% !TeX root = ../../Book2.tex
\OptionalSection{Fourier 级数的深层收敛现象}{b2:sec:0207}

跳跃附近的 Fourier 部分和会发生局部过冲，Dirichlet 核绝对质量的增长会使算子范数增大；这些现象揭示了逐点逼近与整体逼近的差别。另一方面，几乎处处收敛还可以在更广的函数类中成立。比较这些结论，能看清局部正则性、积分范数与收敛方式各自的作用。

\subsection{Gibbs 现象}
\label{b2:subsec:020701}

仍取 \(f(x)=x-1/2\) 于 \((0,1)\)。固定 \(c>0\)，在距跳跃约 \(1/N\) 的点 \(x_N=c/(2\pi N)\) 处，
\[
 S_Nf(x_N)=-\dfrac1\pi\sum_{k=1}^N\dfrac{\sin(kc/N)}k
 \longrightarrow-\dfrac1\pi\int_0^c\dfrac{\sin u}{u}\,\mathrm{d}u.
\]
这是连续函数 \(\sin u/u\)（原点补为 \(1\)）的 Riemann 和。另一方面，\(f(x_N)\to-1/2\)。当 \(c=\pi\) 时，右侧约为 \(-0.58949\)，超出了右侧函数值的极限。这种过冲称为\term[gibbs-xianxiang]{Gibbs 现象}[Gibbs phenomenon]。

\ProseParagraphBreak

为看清过冲确实为正，先注意 \(\int_0^\infty\sin u/u\,\mathrm{d}u=\pi/2\)。可在积分中加入 \(e^{-\varepsilon u}\)，对频率参数求导后积分得到 \(\arctan(1/\varepsilon)\)，再由分部积分的一致尾部界令 \(\varepsilon\downarrow0\)。从 \(\pi\) 开始，正负半波的绝对积分逐段减小，第一段为负，故剩余尾积分严格为负。因此 \(\pi^{-1}\int_0^\pi\sin u/u\,\mathrm{d}u-1/2>0\)。过冲的高度不消失，发生过冲的区域却不断收缩；这与平方误差趋零并不矛盾。

\ProseParagraphBreak

这一计算对应于\cref{b2:fig:02-gibbs}。图中的两种截断使用同一个坐标尺度，可同时比较过冲的高度与区域宽度。
% 本书原创：按第三章锯齿函数部分和精确公式生成。
\begin{figure}[htbp]
\centering
\begin{tikzpicture}[x=1cm,y=1cm,>=stealth,every node/.style={font=\small},line width=.65pt]
\begin{scope}[xshift=0.0cm]
\draw[->,gray] (-.1,2.275)--(5.6,2.275) node[right]{\(x\)};
\draw[->,gray] (2.64,-.05)--(2.64,4.7);
\node at (2.64,5.0) {\(N=8\)};
\draw[gray] (2.57,0.5250)--(2.71,0.5250);\node[left] at (-0.12,0.5250) {\(-1/2\)};
\draw[gray] (2.57,2.2750)--(2.71,2.2750);\node[left] at (-0.12,2.2750) {\(0\)};
\draw[gray] (2.57,4.0250)--(2.71,4.0250);\node[left] at (-0.12,4.0250) {\(1/2\)};
\draw[gray,dashed,thick] plot coordinates {(0.00000,3.60500) (0.05387,3.61357) (0.10775,3.62214) (0.16162,3.63071) (0.21549,3.63928) (0.26937,3.64785) (0.32324,3.65642) (0.37711,3.66500) (0.43098,3.67357) (0.48486,3.68214) (0.53873,3.69071) (0.59260,3.69928) (0.64648,3.70785) (0.70035,3.71642) (0.75422,3.72499) (0.80810,3.73356) (0.86197,3.74213) (0.91584,3.75070) (0.96972,3.75927) (1.02359,3.76784) (1.07746,3.77641) (1.13133,3.78498) (1.18521,3.79356) (1.23908,3.80213) (1.29295,3.81070) (1.34683,3.81927) (1.40070,3.82784) (1.45457,3.83641) (1.50845,3.84498) (1.56232,3.85355) (1.61619,3.86212) (1.67006,3.87069) (1.72394,3.87926) (1.77781,3.88783) (1.83168,3.89640) (1.88556,3.90498) (1.93943,3.91355) (1.99330,3.92212) (2.04718,3.93069) (2.10105,3.93926) (2.15492,3.94783) (2.20880,3.95640) (2.26267,3.96497) (2.31654,3.97354) (2.37041,3.98211) (2.42429,3.99068) (2.47816,3.99925) (2.53203,4.00782) (2.58591,4.01639) (2.63978,4.02496)};
\draw[gray,dashed,thick] plot coordinates {(2.64022,0.52504) (2.69409,0.53361) (2.74797,0.54218) (2.80184,0.55075) (2.85571,0.55932) (2.90959,0.56789) (2.96346,0.57646) (3.01733,0.58503) (3.07120,0.59360) (3.12508,0.60217) (3.17895,0.61074) (3.23282,0.61931) (3.28670,0.62788) (3.34057,0.63645) (3.39444,0.64502) (3.44832,0.65360) (3.50219,0.66217) (3.55606,0.67074) (3.60994,0.67931) (3.66381,0.68788) (3.71768,0.69645) (3.77155,0.70502) (3.82543,0.71359) (3.87930,0.72216) (3.93317,0.73073) (3.98705,0.73930) (4.04092,0.74787) (4.09479,0.75644) (4.14867,0.76502) (4.20254,0.77359) (4.25641,0.78216) (4.31028,0.79073) (4.36416,0.79930) (4.41803,0.80787) (4.47190,0.81644) (4.52578,0.82501) (4.57965,0.83358) (4.63352,0.84215) (4.68740,0.85072) (4.74127,0.85929) (4.79514,0.86786) (4.84902,0.87643) (4.90289,0.88501) (4.95676,0.89358) (5.01063,0.90215) (5.06451,0.91072) (5.11838,0.91929) (5.17225,0.92786) (5.22613,0.93643) (5.28000,0.94500)};
\draw[AnalysisBlue,thick] plot coordinates {(0.00000,3.43226) (0.02039,3.43461) (0.04077,3.43740) (0.06116,3.44063) (0.08154,3.44430) (0.10193,3.44843) (0.12232,3.45301) (0.14270,3.45806) (0.16309,3.46357) (0.18347,3.46954) (0.20386,3.47598) (0.22425,3.48289) (0.24463,3.49027) (0.26502,3.49811) (0.28541,3.50640) (0.30579,3.51515) (0.32618,3.52435) (0.34656,3.53399) (0.36695,3.54407) (0.38734,3.55456) (0.40772,3.56548) (0.42811,3.57679) (0.44849,3.58849) (0.46888,3.60057) (0.48927,3.61300) (0.50965,3.62579) (0.53004,3.63890) (0.55042,3.65231) (0.57081,3.66602) (0.59120,3.68000) (0.61158,3.69423) (0.63197,3.70868) (0.65236,3.72334) (0.67274,3.73818) (0.69313,3.75317) (0.71351,3.76829) (0.73390,3.78351) (0.75429,3.79881) (0.77467,3.81415) (0.79506,3.82951) (0.81544,3.84486) (0.83583,3.86016) (0.85622,3.87539) (0.87660,3.89052) (0.89699,3.90552) (0.91737,3.92034) (0.93776,3.93497) (0.95815,3.94936) (0.97853,3.96348) (0.99892,3.97731) (1.01931,3.99080) (1.03969,4.00393) (1.06008,4.01666) (1.08046,4.02895) (1.10085,4.04077) (1.12124,4.05209) (1.14162,4.06288) (1.16201,4.07310) (1.18239,4.08272) (1.20278,4.09171) (1.22317,4.10004) (1.24355,4.10766) (1.26394,4.11456) (1.28432,4.12071) (1.30471,4.12606) (1.32510,4.13060) (1.34548,4.13429) (1.36587,4.13711) (1.38625,4.13903) (1.40664,4.14002) (1.42703,4.14007) (1.44741,4.13914) (1.46780,4.13721) (1.48819,4.13426) (1.50857,4.13027) (1.52896,4.12523) (1.54934,4.11910) (1.56973,4.11188) (1.59012,4.10355) (1.61050,4.09410) (1.63089,4.08351) (1.65127,4.07176) (1.67166,4.05886) (1.69205,4.04479) (1.71243,4.02954) (1.73282,4.01312) (1.75320,3.99550) (1.77359,3.97670) (1.79398,3.95670) (1.81436,3.93552) (1.83475,3.91315) (1.85514,3.88959) (1.87552,3.86485) (1.89591,3.83895) (1.91629,3.81187) (1.93668,3.78364) (1.95707,3.75427) (1.97745,3.72377) (1.99784,3.69215) (2.01822,3.65943) (2.03861,3.62563) (2.05900,3.59077) (2.07938,3.55486) (2.09977,3.51793) (2.12015,3.48000) (2.14054,3.44110) (2.16093,3.40125) (2.18131,3.36049) (2.20170,3.31883) (2.22208,3.27632) (2.24247,3.23297) (2.26286,3.18883) (2.28324,3.14392) (2.30363,3.09829) (2.32402,3.05197) (2.34440,3.00499) (2.36479,2.95739) (2.38517,2.90921) (2.40556,2.86050) (2.42595,2.81128) (2.44633,2.76161) (2.46672,2.71151) (2.48710,2.66105) (2.50749,2.61025) (2.52788,2.55916) (2.54826,2.50783) (2.56865,2.45630) (2.58903,2.40461) (2.60942,2.35281) (2.62981,2.30095) (2.65019,2.24905) (2.67058,2.19719) (2.69097,2.14539) (2.71135,2.09370) (2.73174,2.04217) (2.75212,1.99084) (2.77251,1.93975) (2.79290,1.88895) (2.81328,1.83849) (2.83367,1.78839) (2.85405,1.73872) (2.87444,1.68950) (2.89483,1.64079) (2.91521,1.59261) (2.93560,1.54501) (2.95598,1.49803) (2.97637,1.45171) (2.99676,1.40608) (3.01714,1.36117) (3.03753,1.31703) (3.05792,1.27368) (3.07830,1.23117) (3.09869,1.18951) (3.11907,1.14875) (3.13946,1.10890) (3.15985,1.07000) (3.18023,1.03207) (3.20062,0.99514) (3.22100,0.95923) (3.24139,0.92437) (3.26178,0.89057) (3.28216,0.85785) (3.30255,0.82623) (3.32293,0.79573) (3.34332,0.76636) (3.36371,0.73813) (3.38409,0.71105) (3.40448,0.68515) (3.42486,0.66041) (3.44525,0.63685) (3.46564,0.61448) (3.48602,0.59330) (3.50641,0.57330) (3.52680,0.55450) (3.54718,0.53688) (3.56757,0.52046) (3.58795,0.50521) (3.60834,0.49114) (3.62873,0.47824) (3.64911,0.46649) (3.66950,0.45590) (3.68988,0.44645) (3.71027,0.43812) (3.73066,0.43090) (3.75104,0.42477) (3.77143,0.41973) (3.79181,0.41574) (3.81220,0.41279) (3.83259,0.41086) (3.85297,0.40993) (3.87336,0.40998) (3.89375,0.41097) (3.91413,0.41289) (3.93452,0.41571) (3.95490,0.41940) (3.97529,0.42394) (3.99568,0.42929) (4.01606,0.43544) (4.03645,0.44234) (4.05683,0.44996) (4.07722,0.45829) (4.09761,0.46728) (4.11799,0.47690) (4.13838,0.48712) (4.15876,0.49791) (4.17915,0.50923) (4.19954,0.52105) (4.21992,0.53334) (4.24031,0.54607) (4.26069,0.55920) (4.28108,0.57269) (4.30147,0.58652) (4.32185,0.60064) (4.34224,0.61503) (4.36263,0.62966) (4.38301,0.64448) (4.40340,0.65948) (4.42378,0.67461) (4.44417,0.68984) (4.46456,0.70514) (4.48494,0.72049) (4.50533,0.73585) (4.52571,0.75119) (4.54610,0.76649) (4.56649,0.78171) (4.58687,0.79683) (4.60726,0.81182) (4.62764,0.82666) (4.64803,0.84132) (4.66842,0.85577) (4.68880,0.87000) (4.70919,0.88398) (4.72958,0.89769) (4.74996,0.91110) (4.77035,0.92421) (4.79073,0.93700) (4.81112,0.94943) (4.83151,0.96151) (4.85189,0.97321) (4.87228,0.98452) (4.89266,0.99544) (4.91305,1.00593) (4.93344,1.01601) (4.95382,1.02565) (4.97421,1.03485) (4.99459,1.04360) (5.01498,1.05189) (5.03537,1.05973) (5.05575,1.06711) (5.07614,1.07402) (5.09653,1.08046) (5.11691,1.08643) (5.13730,1.09194) (5.15768,1.09699) (5.17807,1.10157) (5.19846,1.10570) (5.21884,1.10937) (5.23923,1.11260) (5.25961,1.11539) (5.28000,1.11774)};
\end{scope}
\begin{scope}[xshift=6.6cm]
\draw[->,gray] (-.1,2.275)--(5.6,2.275) node[right]{\(x\)};
\draw[->,gray] (2.64,-.05)--(2.64,4.7);
\node at (2.64,5.0) {\(N=32\)};
\draw[gray] (2.57,0.5250)--(2.71,0.5250);\node[left] at (-0.12,0.5250) {\(-1/2\)};
\draw[gray] (2.57,2.2750)--(2.71,2.2750);\node[left] at (-0.12,2.2750) {\(0\)};
\draw[gray] (2.57,4.0250)--(2.71,4.0250);\node[left] at (-0.12,4.0250) {\(1/2\)};
\draw[gray,dashed,thick] plot coordinates {(0.00000,3.60500) (0.05387,3.61357) (0.10775,3.62214) (0.16162,3.63071) (0.21549,3.63928) (0.26937,3.64785) (0.32324,3.65642) (0.37711,3.66500) (0.43098,3.67357) (0.48486,3.68214) (0.53873,3.69071) (0.59260,3.69928) (0.64648,3.70785) (0.70035,3.71642) (0.75422,3.72499) (0.80810,3.73356) (0.86197,3.74213) (0.91584,3.75070) (0.96972,3.75927) (1.02359,3.76784) (1.07746,3.77641) (1.13133,3.78498) (1.18521,3.79356) (1.23908,3.80213) (1.29295,3.81070) (1.34683,3.81927) (1.40070,3.82784) (1.45457,3.83641) (1.50845,3.84498) (1.56232,3.85355) (1.61619,3.86212) (1.67006,3.87069) (1.72394,3.87926) (1.77781,3.88783) (1.83168,3.89640) (1.88556,3.90498) (1.93943,3.91355) (1.99330,3.92212) (2.04718,3.93069) (2.10105,3.93926) (2.15492,3.94783) (2.20880,3.95640) (2.26267,3.96497) (2.31654,3.97354) (2.37041,3.98211) (2.42429,3.99068) (2.47816,3.99925) (2.53203,4.00782) (2.58591,4.01639) (2.63978,4.02496)};
\draw[gray,dashed,thick] plot coordinates {(2.64022,0.52504) (2.69409,0.53361) (2.74797,0.54218) (2.80184,0.55075) (2.85571,0.55932) (2.90959,0.56789) (2.96346,0.57646) (3.01733,0.58503) (3.07120,0.59360) (3.12508,0.60217) (3.17895,0.61074) (3.23282,0.61931) (3.28670,0.62788) (3.34057,0.63645) (3.39444,0.64502) (3.44832,0.65360) (3.50219,0.66217) (3.55606,0.67074) (3.60994,0.67931) (3.66381,0.68788) (3.71768,0.69645) (3.77155,0.70502) (3.82543,0.71359) (3.87930,0.72216) (3.93317,0.73073) (3.98705,0.73930) (4.04092,0.74787) (4.09479,0.75644) (4.14867,0.76502) (4.20254,0.77359) (4.25641,0.78216) (4.31028,0.79073) (4.36416,0.79930) (4.41803,0.80787) (4.47190,0.81644) (4.52578,0.82501) (4.57965,0.83358) (4.63352,0.84215) (4.68740,0.85072) (4.74127,0.85929) (4.79514,0.86786) (4.84902,0.87643) (4.90289,0.88501) (4.95676,0.89358) (5.01063,0.90215) (5.06451,0.91072) (5.11838,0.91929) (5.17225,0.92786) (5.22613,0.93643) (5.28000,0.94500)};
\draw[AnalysisBlue,thick] plot coordinates {(0.00000,3.56851) (0.02039,3.57759) (0.04077,3.58787) (0.06116,3.59911) (0.08154,3.61103) (0.10193,3.62333) (0.12232,3.63570) (0.14270,3.64780) (0.16309,3.65932) (0.18347,3.66994) (0.20386,3.67939) (0.22425,3.68743) (0.24463,3.69385) (0.26502,3.69852) (0.28541,3.70136) (0.30579,3.70235) (0.32618,3.70155) (0.34656,3.69905) (0.36695,3.69506) (0.38734,3.68979) (0.40772,3.68353) (0.42811,3.67661) (0.44849,3.66939) (0.46888,3.66222) (0.48927,3.65549) (0.50965,3.64958) (0.53004,3.64481) (0.55042,3.64150) (0.57081,3.63992) (0.59120,3.64027) (0.61158,3.64269) (0.63197,3.64726) (0.65236,3.65397) (0.67274,3.66273) (0.69313,3.67340) (0.71351,3.68573) (0.73390,3.69945) (0.75429,3.71419) (0.77467,3.72957) (0.79506,3.74516) (0.81544,3.76052) (0.83583,3.77521) (0.85622,3.78880) (0.87660,3.80089) (0.89699,3.81113) (0.91737,3.81923) (0.93776,3.82496) (0.95815,3.82816) (0.97853,3.82880) (0.99892,3.82688) (1.01931,3.82255) (1.03969,3.81601) (1.06008,3.80757) (1.08046,3.79758) (1.10085,3.78650) (1.12124,3.77482) (1.14162,3.76305) (1.16201,3.75174) (1.18239,3.74142) (1.20278,3.73262) (1.22317,3.72582) (1.24355,3.72142) (1.26394,3.71979) (1.28432,3.72118) (1.30471,3.72574) (1.32510,3.73353) (1.34548,3.74447) (1.36587,3.75840) (1.38625,3.77501) (1.40664,3.79392) (1.42703,3.81462) (1.44741,3.83654) (1.46780,3.85905) (1.48819,3.88147) (1.50857,3.90310) (1.52896,3.92323) (1.54934,3.94122) (1.56973,3.95642) (1.59012,3.96832) (1.61050,3.97645) (1.63089,3.98050) (1.65127,3.98027) (1.67166,3.97570) (1.69205,3.96691) (1.71243,3.95414) (1.73282,3.93781) (1.75320,3.91847) (1.77359,3.89680) (1.79398,3.87360) (1.81436,3.84978) (1.83475,3.82627) (1.85514,3.80409) (1.87552,3.78423) (1.89591,3.76766) (1.91629,3.75529) (1.93668,3.74794) (1.95707,3.74629) (1.97745,3.75086) (1.99784,3.76199) (2.01822,3.77981) (2.03861,3.80424) (2.05900,3.83493) (2.07938,3.87133) (2.09977,3.91263) (2.12015,3.95779) (2.14054,4.00557) (2.16093,4.05453) (2.18131,4.10306) (2.20170,4.14945) (2.22208,4.19186) (2.24247,4.22844) (2.26286,4.25731) (2.28324,4.27667) (2.30363,4.28477) (2.32402,4.28003) (2.34440,4.26104) (2.36479,4.22661) (2.38517,4.17581) (2.40556,4.10801) (2.42595,4.02288) (2.44633,3.92042) (2.46672,3.80097) (2.48710,3.66521) (2.50749,3.51414) (2.52788,3.34909) (2.54826,3.17167) (2.56865,2.98375) (2.58903,2.78741) (2.60942,2.58494) (2.62981,2.37873) (2.65019,2.17127) (2.67058,1.96506) (2.69097,1.76259) (2.71135,1.56625) (2.73174,1.37833) (2.75212,1.20091) (2.77251,1.03586) (2.79290,0.88479) (2.81328,0.74903) (2.83367,0.62958) (2.85405,0.52712) (2.87444,0.44199) (2.89483,0.37419) (2.91521,0.32339) (2.93560,0.28896) (2.95598,0.26997) (2.97637,0.26523) (2.99676,0.27333) (3.01714,0.29269) (3.03753,0.32156) (3.05792,0.35814) (3.07830,0.40055) (3.09869,0.44694) (3.11907,0.49547) (3.13946,0.54443) (3.15985,0.59221) (3.18023,0.63737) (3.20062,0.67867) (3.22100,0.71507) (3.24139,0.74576) (3.26178,0.77019) (3.28216,0.78801) (3.30255,0.79914) (3.32293,0.80371) (3.34332,0.80206) (3.36371,0.79471) (3.38409,0.78234) (3.40448,0.76577) (3.42486,0.74591) (3.44525,0.72373) (3.46564,0.70022) (3.48602,0.67640) (3.50641,0.65320) (3.52680,0.63153) (3.54718,0.61219) (3.56757,0.59586) (3.58795,0.58309) (3.60834,0.57430) (3.62873,0.56973) (3.64911,0.56950) (3.66950,0.57355) (3.68988,0.58168) (3.71027,0.59358) (3.73066,0.60878) (3.75104,0.62677) (3.77143,0.64690) (3.79181,0.66853) (3.81220,0.69095) (3.83259,0.71346) (3.85297,0.73538) (3.87336,0.75608) (3.89375,0.77499) (3.91413,0.79160) (3.93452,0.80553) (3.95490,0.81647) (3.97529,0.82426) (3.99568,0.82882) (4.01606,0.83021) (4.03645,0.82858) (4.05683,0.82418) (4.07722,0.81738) (4.09761,0.80858) (4.11799,0.79826) (4.13838,0.78695) (4.15876,0.77518) (4.17915,0.76350) (4.19954,0.75242) (4.21992,0.74243) (4.24031,0.73399) (4.26069,0.72745) (4.28108,0.72312) (4.30147,0.72120) (4.32185,0.72184) (4.34224,0.72504) (4.36263,0.73077) (4.38301,0.73887) (4.40340,0.74911) (4.42378,0.76120) (4.44417,0.77479) (4.46456,0.78948) (4.48494,0.80484) (4.50533,0.82043) (4.52571,0.83581) (4.54610,0.85055) (4.56649,0.86427) (4.58687,0.87660) (4.60726,0.88727) (4.62764,0.89603) (4.64803,0.90274) (4.66842,0.90731) (4.68880,0.90973) (4.70919,0.91008) (4.72958,0.90850) (4.74996,0.90519) (4.77035,0.90042) (4.79073,0.89451) (4.81112,0.88778) (4.83151,0.88061) (4.85189,0.87339) (4.87228,0.86647) (4.89266,0.86021) (4.91305,0.85494) (4.93344,0.85095) (4.95382,0.84845) (4.97421,0.84765) (4.99459,0.84864) (5.01498,0.85148) (5.03537,0.85615) (5.05575,0.86257) (5.07614,0.87061) (5.09653,0.88006) (5.11691,0.89068) (5.13730,0.90220) (5.15768,0.91430) (5.17807,0.92667) (5.19846,0.93897) (5.21884,0.95089) (5.23923,0.96213) (5.25961,0.97241) (5.28000,0.98149)};
\end{scope}
\end{tikzpicture}
\caption[Gibbs 过冲与尺度收缩]{Gibbs 过冲与尺度收缩。灰色虚线为锯齿函数在跳跃两侧的值，蓝线为对称部分和；两图的坐标尺度相同。截断阶数增大后，过冲集中于更窄的区域，其相对跳跃高度仍不趋零。}
\label{b2:fig:02-gibbs}
\end{figure}


\subsection{对称部分和的\texorpdfstring{\(L^1\)}{L¹}收敛障碍}
\label{b2:subsec:020702}

\begin{proposition}\label{b2:prop:02-lebesgue-constants}
存在绝对常数 \(c,C>0\)，使 \(N\geq2\) 时
\[
 c\cdot\log(N+1)\leq\|D_N\|_1\leq C\cdot\log(N+1).
\]
此外，\(\|S_N\|_{L^1\to L^1}=\|D_N\|_1\)。
\end{proposition}
\begin{proof}
在 \(0<t\leq1/2\) 上，\(|D_N(t)|\leq\min\{2N+1,(2t)^{-1}\}\)，分成 \((0,1/(2N+1))\) 和其余部分积分，得到上界。下界利用 \(\sin(\pi t)\leq\pi t\)；令 \(u=(2N+1)\pi t\)，可得
\[
 \|D_N\|_1\geq\dfrac2\pi\int_0^{(2N+1)\pi/2}\dfrac{|\sin u|}{u}\,\mathrm{d}u.
\]
在每个完整区间 \([j\pi,(j+1)\pi]\) 上，以 \(((j+1)\pi)^{-1}\) 下估 \(1/u\)，而 \(\int_{j\pi}^{(j+1)\pi}|\sin u|\,\mathrm{d}u=2\)。求和即得调和级数的对数下界。

卷积和 Tonelli 给出算子范数不超过 \(\|D_N\|_1\)。反向取非负单位积分的窄区间核 \(q_\varepsilon=\varepsilon^{-1}\mathbf1_{(0,\varepsilon)}\)。固定 \(N\) 时，\(q_\varepsilon*D_N\to D_N\) 一致成立，因为 \(D_N\) 连续；且 \(\|q_\varepsilon\|_1=1\)。故算子范数至少为 \(\|D_N\|_1\)。
\end{proof}

\begin{corollary}\label{b2:cor:02-l1-failure}
存在 \(f\in L^1(\mathbb T)\)，使 \(\sup_N\|S_Nf\|_1=\infty\)。因而其 Fourier 部分和不在 \(L^1\) 中收敛。
\end{corollary}
\begin{proof}
若每个 \(f\) 的部分和范数都有界，由第一卷定理 8.3.3（一致有界原理），算子族 \((S_N)\) 的范数也应一致有界，这与前一命题矛盾。
\end{proof}

同一个函数的 Fejér 平均仍按 \(L^1\) 范数收敛。失败属于所选求和方法，不是 Fourier 系数失去了唯一确定函数的能力。

\subsection{发散例子}
\label{b2:subsec:020703}

\begin{proposition}\label{b2:prop:02-continuous-divergence}
对任意预先指定的 \(x_0\in\mathbb T\)，存在连续周期函数 \(f\)，使 \(\sup_N|S_Nf(x_0)|=\infty\)。
\end{proposition}
\begin{proof}
将 \(T_Nf=S_Nf(x_0)\) 看作 \(C(\mathbb T)\) 上的连续线性泛函。积分公式给出 \(\|T_N\|\leq\|D_N\|_1\)。为证明反向不等式，取连续函数
\[
 f_\varepsilon(x_0-t)=\dfrac{D_N(t)}{|D_N(t)|+\varepsilon}.
\]
它的上确界范数至多为 \(1\)，且 \(T_Nf_\varepsilon=\int D_N^2/(|D_N|+\varepsilon)\to\|D_N\|_1\)，其中控制函数为 \(|D_N|\)。故 \(\|T_N\|=\|D_N\|_1\to\infty\)。再用第一卷定理 8.3.3，存在固定的 \(f\) 使 \((T_Nf)\) 无界。
\end{proof}

这个存在性证明没有给出一串闭式系数，但确实产生了一个对所有 \(N\) 固定的连续函数。只为每个 \(N\) 另选一个坏函数，尚不能证明存在单个发散例子；一致有界原理在这里承担的正是这一量词转换。

\subsection{Carleson 定理概览}
\label{b2:subsec:020704}

\begin{claim}[Carleson，深层收敛定理]\label{b2:claim:02-carleson}
取 \(\mathbb T=\mathbb R/\mathbb Z\) 的归一化 Lebesgue 测度。若 \(f\in L^2(\mathbb T;\mathbb C)\)，令 \(\widehat f(k)=\int_{\mathbb T}f(t)e^{-2\pi\mathrm i kt}\,\mathrm dt\)、\(S_Nf(x)=\sum_{|k|\le N}\widehat f(k)e^{2\pi\mathrm i kx}\)，\(N\in\mathbb N_0\)。则对称 Fourier 部分和 \(S_Nf(x)\) 在几乎处处收敛到 \(f(x)\)。
\end{claim}

这一\term[carleson-dingli]{Carleson 定理}[Carleson's theorem]的证明超出了本节的范围。Sheldon Axler 的《Measure, Integration \& Real Analysis》\cite[p. 352, Section 11A]{Axler2020Measure}介绍了该结论与连续函数发散例子的关系。其困难在于同时控制全部截断频率；固定一个 \(N\) 的平方范数估计，以及本章抽子列的办法，都不能代替这种控制。

\ProseParagraphBreak

“几乎处处”允许存在例外零测集，所以定理与在指定点发散的连续函数不冲突。对于一般的 \(L^1(\mathbb T)\) 函数，不能保证对称 Fourier 部分和几乎处处收敛。Kolmogorov 构造了一个 \(f\in L^1(\mathbb T)\)，使对称部分和序列 \((S_Nf(x))_{N\geq0}\) 在几乎处处发散；上述资料同页介绍了这一结果。这是存在性反例，并非每个可积函数都发散：例如三角多项式的部分和在充分大时就等于它自身。

\ProseParagraphBreak

Kolmogorov 的具体构造在此略去。这个反例涉及几乎处处发散，与前面讨论的 \(L^1\) 范数收敛失败是两个不同结论。读者首先应掌握普通部分和、正核平均、范数收敛与点态收敛之间的区别，再进入深层收敛理论。
